\documentclass[a4paper,oneside,10pt]{article}

\usepackage{fontspec}
\setmainfont{CMU Serif Roman}
\usepackage{libertine}

\usepackage[greek,english]{babel}

\usepackage[colorlinks = true,
            linkcolor = blue,
            urlcolor  = blue,
            citecolor = blue,
            anchorcolor = blue]{hyperref}

\useshorthands{;}
\defineshorthand{;}{?}

\usepackage{graphicx,epstopdf}
\usepackage{caption}

\usepackage{float}
\usepackage{listings}
\usepackage{parskip}
\usepackage{fancyhdr}
\usepackage[export]{adjustbox}
\setcounter{secnumdepth}{5}
\setcounter{tocdepth}{5}

\pagestyle{fancy}
\fancyhf{}
\lfoot{}
\rfoot{\thepage}

\renewcommand{\headrulewidth}{0.5pt}
\renewcommand{\footrulewidth}{0.5pt}

\usepackage{tikz}
\newcommand{\roundpic}[4][]{
  \tikz\node [circle, minimum width = #2,
    path picture = {
      \node [#1] at (path picture bounding box.center) {
        \includegraphics[width=#3]{#4}};
    }] {};}

\usepackage{multicol}

\pagestyle{fancy}
\fancyhf{}
\rhead{Curriculum Vitae}
\lhead{Iosif Angelidis}

\begin{document}

\begin{multicols}{2} 
\roundpic{5cm}{5cm}{IMG_20190204_195706_ex.jpg}
\columnbreak

\textlatin{Name}: \textlatin{Iosif}

\textlatin{Surname}: \textlatin{Angelidis}

\textlatin{Birth}: \textlatin{June 16, 1993}

\textlatin{Location}: \textlatin{Athens, Greece}

\textlatin{Position}: \textlatin{IT professional}

\textlatin{Mobile}: \textlatin{\href{tel:306947747843}{+30 694 7747 843}}

\textlatin{e-mail}: \textlatin{\href{mailto:iosif.angelidis@gmail.com}{iosif.angelidis@gmail.com}}


\end{multicols}

\subsection*{Contact}
\addcontentsline{toc}{subsection}{Contact}

\begin{multicols}{10}
\begin{itemize}

\item[]\href{https://iosang.github.io}{\includegraphics[scale=0.08]{assets/home.pdf}}

\item[]\href{https://www.facebook.com/metimdjai}{\includegraphics[scale=0.2]{assets/facebook-original.pdf}}

\item[]\href{https://www.linkedin.com/in/iosif-angelidis/}{\includegraphics[scale=0.2]{assets/linkedin-original.pdf}}

\item[]\href{https://twitter.com/metimdjai}{\includegraphics[scale=0.2]{assets/twitter-original.pdf}}

\item[]\href{https://github.com/iosang}{\includegraphics[scale=0.2]{assets/github-original.pdf}}

\item[]\href{https://gitlab.com/metimdjai}{\includegraphics[scale=0.2]{assets/gitlab-original.pdf}}

\item[]\href{https://bitbucket.org/Metimdjai/}{\includegraphics[scale=0.2]{assets/bitbucket-original.pdf}}

\item[]\href{https://dblp.uni-trier.de/pers/hd/a/Angelidis:Iosif}{\includegraphics[scale=0.15]{assets/dblp.pdf}}

\item[]\href{https://scholar.google.gr/citations?user=r8jJLocAAAAJ&hl=en}{\includegraphics[scale=0.3]{assets/icons8-google-scholar.pdf}}

\item[]\href{skype:joeguru1?chat}{\includegraphics[scale=0.09]{assets/skype.pdf}}

\end{itemize}
\end{multicols}

\begin{figure}[hbt!]
	\centering
	\includegraphics[scale=0.315]{qrcode.png}%
\end{figure}

\subsection*{Short summary}
\addcontentsline{toc}{subsection}{Short summary}

\begin{sloppypar}
	Iosif Angelidis is an IT professional with four (4) years of experience in software development stemming. His IT skills spam from IT programming with actual experience in C, C++, Java and Python, among other languages. This knowledge is gained both from self-motive to invest in new skill-sets and as part of his two years engagement in a Named Entity Recognizer thesis, which recognizes and extracts persons, organizations, geo-political entities, legal references, public documents and unofficial geographical landmarks found in the Government Gazette by utilizing Neural Networks and Deep Learning technologies.

	During his military tour, Iosif has been occupied in the IT Division offering his experience in supporting 3 projects: MIP 3, ΣΧΚ and ΣΔΕΠ. This support came in the form of major refactoring, implementation of new features, optimization and critical performance gains, as well as general user support. His responsibilities further included the training of new staff and knowledge transferring.

	Working at UBITECH since July 2020, he has contributed technical implementations, partially managed and delivered reports for H2020 projects \href{https://cordis.europa.eu/project/id/780938}{ChildRescue}, \href{https://cordis.europa.eu/project/id/760792}{FENIX}, \href{https://cordis.europa.eu/project/id/882897}{Search And Rescue}, \href{https://cordis.europa.eu/project/id/871703}{TheFSM}, \href{https://cordis.europa.eu/project/id/959879}{GLASS}, the EU accelerator DigiCirc project \href{https://digicirc.eu/irecirc/}{iReCirc } and ΕΣΠΑ projects \href{https://endorse.biosim.ntua.gr/}{ENDORSE }, 3DFollicleAI and PHAEDON. While improving his technical skills, he self-invests in learning management and software architect skills as well, resulting in him participating in the partial management of some of the aforementioned projects, as well as conducting requirement analysis, proposing conceptual, technical and deployment architectures and applying Agile methodology and User stories in the analysis process. Additionally, his contributions include deliverables writing, thorough technical documentations and demonstrators.

	Iosif Angelidis has also shown particular interest in Artificial Intelligence courses both as a tutor (assistant professor) and as a practitioner. 
	Last, but not least, Iosif has become a member of various communities among which the most notable ones pertain to the Management of Data and Information Knowledge Group (\href{http://www.madgik.di.uoa.gr}{MADgIK}) and the \href{http://ai.di.uoa.gr/}{AI group} of the University of Athens. 

	\textbf{Military obligations fulfilled.}
	
	
\end{sloppypar}

\subsection*{Skillsets}
\addcontentsline{toc}{subsection}{Skillsets}

	\begin{itemize}

		\item \textbf{Soft}: Requirement analysis (functional, non-functional, technical, data, user requirements), Agile methodology/User stories, Conceptual/Technical/Deployment architecture design, presentation (slides, sequence diagrams, activity diagrams etc.), communication, team player, self-driven

		\item \textbf{Programming Languages, Applications, Security, Access control, DevOps \& Frameworks}: C, C++, Java, Python, PHP, Javascript, REST, SOAP, Android Studio development, TLS/SSL protocols, MQTT protocol, bower, npm, HTML5, CSS3, XML, RDF, OWL, Ajax, Gradle, Maven, EJBs, JSF, PrimeFaces, ORM (Hibernate), RSA encryption, Symmetric encryption, Hybrid encryption, ACL, RBAC (with/without domains), ABAC (with/without domains), PERM meta-model, API Gateway, Docker, CI/CD, Unit testing, OpenGL, Bootstrap, Sweet alert, Bootstrap Validator, jQuery, Leaflet JS, Particle JS, Datatables, Selectize, Querybuilder, jquery-validator, jquery-steps, papa-parse, Canvas JS, gson, jCasbin, Spring ecosystem (Boot, Security, Data, Thymeleaf, etc.), Swagger documentation, Keras, TF, SciKit Learn, Numpy, Pandas, flask.

		\item \textbf{Infrastructure} 

			\begin{itemize}
				\item \textbf{Operating Systems}: GNU/LINUX (including advanced distros like Arch and Manjaro), Microsoft Windows, MacOS 

				\item \textbf{Shell Scripting}: bash, fish, zsh, tcsh, csh

				\item \textbf{Server Management \& Databases}: Apache Tomcat, Wildfly/Jboss, Glassfish, Apache HTTP Server/XAMPP, PHPMyAdmin, MySQL Workbench, pgAdmin, SQL, PostgreSQL/PostGIS, Oracle 11g/12c, MongoDB, Elasticsearch, Redis, (Geo/st)SPARQL, SQL Server Management Studio, Toad

				\item \textbf{Integrated Desktop Environments}: Intellij, PyCharm, WebStorm, PHPStorm, Vistual Studio Code, Netbeans, Jupyter notebooks

			\end{itemize}

		\item \textbf{Geospatial Systems}: Sextant, Strabon, QGis

		\item \textbf{Version Control Systems}: Mercurial, Git, CVS, GitHub, GitLab, Bitbucket

		\item \textbf{Other}: AOP, Design patterns, \LaTeX, Microsoft Office, Matlab, Gimp

	\end{itemize}

\subsection*{Working Experience}
\addcontentsline{toc}{subsection}{Working Experience}

	\begin{itemize}

		\item \textbf{2020 - present UBITECH Intelligent Computing Research Associate}
		
			During his time at UBITECH, he has been tasked with partial management of H2020 projects: \href{https://cordis.europa.eu/project/id/882897}{Search And Rescue}, \href{https://cordis.europa.eu/project/id/871703}{TheFSM}, \href{https://cordis.europa.eu/project/id/959879}{GLASS}, the EU accelerator DigiCirc project \href{https://digicirc.eu/irecirc/}{iReCirc } and ΕΣΠΑ projects \href{https://endorse.biosim.ntua.gr/}{ENDORSE }, 3DFollicleAI and PHAEDON. His responsibilities additionally include frequent coordination telcos with partners to meet the projects' needs. Additionally, he conducts requirement analysis on functional, non-functional, technical, data, user requirements based on Agile methodology and User stories. Moreover, he is tasked with technical implemementation, API documentation and relevant responsibilities on the above projects (as well as \href{https://cordis.europa.eu/project/id/780938}{ChildRescue}, \href{https://cordis.europa.eu/project/id/760792}{FENIX} and 3DFollicleAI). Finally, he is also tasked with DevOps contributions such as dockerization, deployment and CI/CD pipelines, as well as contributions to deliverables writing, technical documentation and demonstrators.

            Programming skillsets developed during this time involve Java, HTML 5, CSS 3, JavaScript, Leaflet JS, REST, IntelliJ, Android Stutio, Keras, Tensorflow, Pandas, Scikit-learn, ABAC, RBAC, ACL, Elasticsearch, MongoDB, MySQL, PHPmyAdmin, Maven, Pycharm, Docker, CI/CD.
			
			Management skills developed during this time involve requirement analysis (functional, non-functional, technical, data, user requirements), Agile methodology, User stories, conceptual/technical/deployment architecture design.

		\item \textbf{2019 - 2020 Military Service IT Division Software Developer}

			During his military tour, he was tasked with offering his experience in supporting 3 projects: the implementation of the Multi-Interoperability Protocol 3.0 (MIP 3), the Crisis Management System (ΣΧΚ) and the Information Management-Control System (ΣΔΕΠ). These projects required major code refactoring, optimizations, critical performance gains (e.g., database requests taking about 3 seconds compared to 15 minutes) and implementation of brand new features, involving classified data as well as geospatial information. Additionally, he was responsible for training new staff and providing user support for the platforms.

			Programming skillsets developed during this time involve Java, EJBs, JSF, PrimeFaces, Oracle 11g/12c, Hibernate, HTML 5, CSS 3, JavaScript, Leaflet JS, REST, SOAP, XML, SQL Server Management Studio, Toad, WildFly/JBoss, IntelliJ.

		\item \textbf{2017 - 2019 University of Athens IT Programmer}

			In the context of his MSc thesis, ``Named Entity Recognition and Linking in Greek Legislation'', under the supervision of professor Manolis Koubarakis and PhD Candidate Ilias Chalkidis, he developed and maintained a Named Entity Recognizer which recognizes and extracts persons, organizations, geo-political entities, legal references, public documents and unofficial geographical landmarks found in the Government Gazette by utilizing Neural Networks and Deep Learning technologies.

			Programming skillsets developed during this time involve Python, Keras, TF, SciKit Learn, Numpy, Pandas, (Geo/st)SPARQL, REST, HTML 5, CSS 3, RDF, OWL, Maven, PyCharm.

	\end{itemize}

\subsection*{Languages}
\addcontentsline{toc}{subsection}{Languages}

	\begin{itemize}

		\item \textbf{English}: Certificate of Proficiency in English, 2009 University of Michigan, First Certificate in English (Grade B), 2008 University of Cambridge

	\end{itemize}	

\subsection*{Education roadmap}
\addcontentsline{toc}{subsection}{Education roadmap}

	\begin{itemize}

		\item \textbf{2016 - 2018: MSc degree in Information and Data Management}, National and Kapodistrian University of Athens, Department of Informatics and Telecommunications. \textbf{Grade: 9.39/10.0}

		\item \textbf{2011 - 2016: Ptychion (4-year degree with thesis) in Computer Science - Theoretical Informatics}, National and Kapodistrian University of Athens, Department of Informatics and Telecommunications. \textbf{Grade: 9.34/10.0}

		\item \textbf{2007: Technical Diploma ECDL} in Databases, Concepts of IT, Word Processing, Presentations, Information \& Communication, Spreadsheets, Using the Computer and Managing Files

	\end{itemize}

\subsection*{Hobbies and interests}
\addcontentsline{toc}{subsection}{Hobbies and interests}

	\begin{itemize}

		\item Exploring the open source community, learning about new technologies and tweaking computer systems.

		\item Long time gamer, almost all genres. Debating about games and their various aspects (gameplay, story, mechanics, gimmicks etc.). 

		\item Dungeons and Dragons roleplayer. Emphasis on interesting character backgrounds with deep motivations. Currently participating in two campaigns.
		
		\item Aspiring DJ with a twist: steaming mostly music featured in games. \href{https://zeno.fm/the-hermit}{Online station} (Auto DJ streams tunes when offline).

	\end{itemize}

\end{document}
